
%%%%%%%%%%%%%%%%%%%%%%%%%%%%%%%%%%%%%%%%%%%%%%%%%%%%%%%%%%%%%%%%%%%%%%%%%%%%%%

%\section{Sufficiency of Condition~{\MixCond}}
%\label{sec: sufficiency of condition mc}
%\input{5_sufficiency_mc.tex}

In this section we prove that Condition~{\MixCond} in Theorem~\ref{thm: miscibility characterization}(a) 
is sufficient for perfect mixability. A perfect-mixing graph constructed in our argument has
precision at most $1$, showing also the first part of Theorem~\ref{thm: miscibility characterization}(b).

Assume that $n\ge 4$. Let $C$ with $\mu = \average(C)$ and 
$C\cup\braced{\mu}\ingroundset\integers$ be the input configuration,
and assume that $C$ satisfies Condition~{\MixCond}. 
The outline of our proof is as follows:
%
\begin{itemize}
	\item First we prove that $C$ is perfectly mixable with precision $0$ when $n$ is a power of $2$.
		This easily extends to configurations $C$ called \emph{near-final}, 
		which are disjoint unions of multisets with the same average and cardinalities being powers of $2$. 
		In particular, this proves Theorem~\ref{thm: miscibility characterization}(a) for $n=4$.
		
	\item Next, we give a proof for $n\ge 7$. The basic idea of the proof is
		to define an invariant~(I) and show that any configuration that
		satisfies~(I) has a pair of droplets whose mixing either preserves
		invariant~(I) or produces a near-final configuration. 
		Condition~(I) is stronger than~{\MixCond} (it implies~{\MixCond}, but not
		vice versa), but we show that
		that any configuration that satisfies Condition~{\MixCond} can be
		modified to satisfy~(I).
		
	\item We then give separate proofs for $n=5,6$.
		The proof for $n=5$ is similar to the case $n\ge 7$, but
		it requires a more subtle invariant.
		The proof for $n=6$ is derived by minor modifications to the proof for $n=5$.
\end{itemize}
